\section{Background}
\label{sec:background}

This section establishes notation and reviews the four theoretical foundations
that Infon builds on.

% -----------------------------------------------------------------------
\subsection{Dempster--Shafer Theory}
\label{sec:background:ds}

Dempster--Shafer (DS) theory~\cite{dempster1967upper,shafer1976mathematical}
generalises Bayesian probability by allowing belief to be assigned to
\emph{sets} of hypotheses rather than individual hypotheses, thereby
distinguishing \emph{uncertainty} (evidence spread over possibilities) from
\emph{ignorance} (absence of evidence).

\paragraph{Frame of discernment.}
Let $\Omega$ be a finite, mutually exclusive, exhaustive set of hypotheses
called the \emph{frame of discernment}.  Infon uses
\[
  \Omega = \{S,\, R,\, U\},
\]
where $S$ = Supports, $R$ = Refutes, and $U$ = Uncertain/mixed.

\paragraph{Basic probability assignment.}
A \emph{basic probability assignment} (bpa) is a function
$m \colon 2^{\Omega} \to [0,1]$ satisfying
\[
  m(\emptyset) = 0 \qquad \text{and} \qquad \sum_{A \subseteq \Omega} m(A) = 1.
\]
Each focal element $A$ with $m(A) > 0$ represents a portion of belief
committed to ``the truth lies somewhere in $A$'' without further
discrimination.

\paragraph{Vacuous element.}
The \emph{vacuous} bpa commits all mass to the full frame:
\[
  m(\Theta) = 1, \quad m(A) = 0 \text{ for all } A \subsetneq \Omega.
\]
Here $\Theta$ denotes $\Omega$ itself when used as a focal element; the
notation $m(\Theta)$ refers to the mass on this total-ignorance element.
It is the \emph{correct} encoding for ``No Evidence Identified'' (NEI): it
expresses that the available evidence does not constrain the hypothesis at all,
as opposed to $m(U) > 0$, which would assert that conflicting evidence
\emph{was} observed.

\paragraph{Dempster's combination rule.}
Given two independent bpas $m_1$ and $m_2$ over the same frame, their
\emph{orthogonal sum} $m = m_1 \oplus m_2$ is defined for all $A \neq
\emptyset$ as
\[
  m(A) = \frac{1}{1 - K} \sum_{\substack{B \cap C = A \\ B,\,C \subseteq \Omega}} m_1(B)\, m_2(C),
\]
where the normalisation constant $K = \sum_{B \cap C = \emptyset} m_1(B)\, m_2(C)$
accounts for conflicting evidence.  Combination is commutative and associative,
so evidence from multiple infons can be fused incrementally.

\paragraph{Four-part mass in Infon.}
Throughout this paper, a verdict mass is written as the tuple
$(m_S,\, m_R,\, m_U,\, m_\Theta)$, where
\[
  m_S = m(\{S\}), \quad
  m_R = m(\{R\}), \quad
  m_U = m(\{U\}), \quad
  m_\Theta = m(\Theta),
\]
and $m_S + m_R + m_U + m_\Theta = 1$.
An \emph{unsupported} claim has $m_\Theta \approx 1$ (no focal evidence);
a \emph{supported} claim has $m_S \approx 1$; a \emph{refuted} claim has
$m_R \approx 1$.  The chain-mass accumulation and GNN scoring described in
Sections~\ref{sec:architecture} and~\ref{sec:implementation} produce masses
in this form.

% -----------------------------------------------------------------------
\subsection{Situation Semantics and Infons}
\label{sec:background:infons}

Infon's unit of extracted knowledge is the \emph{infon}, drawn from
Barwise and Perry's situation semantics~\cite{barwise1983situations}.

\paragraph{Infon triple.}
An infon is a typed quadruple
\[
  \langle\!\langle \mathit{pred},\; \mathit{subj},\; \mathit{obj};\; \mathit{pol} \rangle\!\rangle,
\]
where $\mathit{pred}$ is a predicate label, $\mathit{subj}$ and $\mathit{obj}$
are entity anchors, and $\mathit{pol} \in \{+1, -1\}$ records whether the
predicate holds ($+1$, affirmed) or is negated ($-1$).

\paragraph{Grounding.}
Each infon is grounded to its source: a \emph{cassette} identifier (the SHA-256
content hash of the document), a byte \emph{offset} into the cassette's record
stream, and an extraction \emph{timestamp}.  This grounding lets Infon cite
the exact passage that generated each mass contribution and supports
time-travel queries over snapshot manifests.

\paragraph{Why infons?}
Infons are the natural unit for two reasons.  First, their polarity field
directly encodes negation, so refuting evidence is represented symmetrically
with supporting evidence.  Second, the triple structure $(\mathit{pred},
\mathit{subj}, \mathit{obj})$ maps cleanly onto a hypergraph edge, enabling
MCTS traversal along predicate-typed chains.

% -----------------------------------------------------------------------
\subsection{SPLADE Sparse Retrieval}
\label{sec:background:splade}

SPLADE~\cite{formal2021splade} (Sparse Lexical and Expansion model) trains a
BERT-family encoder to produce high-dimensional \emph{sparse} activations over
a fixed vocabulary.  For a text sequence $t$, the SPLADE encoder produces a
vector $\mathbf{s} \in \mathbb{R}^{|\mathcal{V}|}$ whose $k$-th component is
\[
  s_k = \log\!\bigl(1 + \mathrm{ReLU}(h_k)\bigr),
\]
where $h_k$ is the aggregated BERT hidden state for vocabulary token $k$.
Only a small fraction of dimensions are non-zero, so $\mathbf{s}$ can be
stored and intersected efficiently with an inverted index.

\paragraph{SPLADE-tiny in Infon.}
Infon uses the \texttt{naver/splade-v3-lexical} distilled variant
(17~MB, Apache~2.0 licence).  It runs on CPU in under 2~seconds cold-start
and requires no GPU or external API.  Each ingested document is encoded into a
sparse anchor vector; at query time, the query anchor is matched against
cassette-level bounding boxes in the manifest pruner, and the top-scoring
cassettes are hydrated for full scoring.  This design keeps retrieval
complexity proportional to the number of \emph{relevant} cassettes rather than
the total corpus size.

% -----------------------------------------------------------------------
\subsection{Sheaf Graph Neural Networks}
\label{sec:background:sheaf_gnn}

A sheaf graph neural network~\cite{bodnar2022sheaf} augments a standard message-passing
GNN with \emph{restriction maps}: for each edge $(u, v)$ in the graph, a
learnable linear map $F_{u \trianglelefteq e} \colon \mathbb{R}^d \to
\mathbb{R}^k$ projects node $u$'s feature vector onto a shared
\emph{edge stalk}.  The sheaf Laplacian $\Delta_F$ is defined as
\[
  \Delta_F = B_F^\top B_F,
\]
where $B_F$ is the coboundary operator constructed from the restriction maps.
Its first cohomology group $H^1$ captures the degree to which adjacent node
representations are \emph{inconsistent} with the restriction maps --- a
natural anomaly signal for relational data.

\paragraph{Use in Infon.}
Infon's 140k-parameter sheaf GNN (\texttt{SheafHypergraphEncoder}) uses
per-relation-kind restriction maps, so different predicate types carry
different geometric transformations.  The encoder's chain-verdict head
produces a logit over the mass components $(m_S, m_R, m_U, m_\Theta)$ that serves as the MCTS terminal
scorer.  The $H^1$ discrepancy is separately exposed as an anomaly signal on
reportive edges.  The GNN is trained once on the synthetic corpus
(\S\ref{sec:implementation}) with known ground-truth verdicts, then frozen and
shipped with the package.  At inference time it adds no training overhead; it
functions as a learned heuristic within the MCTS search.
